% =============================================================================
% Section 8: Yocto Build System for Raspberry Pi 5
% =============================================================================

\section{Yocto Build System
}
\label{sec:yocto-build}

This section documents the Yocto-based build system for the STAR robot's Raspberry Pi 5 controller.

\subsection{Overview}

The STAR robot uses the Yocto Project (Scarthgap release) with OpenEmbedded to build a custom Linux distribution for the Raspberry Pi 5. This approach provides:

\begin{itemize}
    \item Full control over the root filesystem
    \item ROS2 Jazzy integration via meta-ros
    \item Custom recipes for STAR-specific packages
    \item Reproducible builds
\end{itemize}

    \subsection{Layer Structure}

    The build uses the following Yocto layers :

\begin{table}[h]
\centering
\begin{tabular}{|l|l|}
\hline
\textbf{Layer} & \textbf{Purpose} \\
\hline
poky/meta & Core OpenEmbedded layer \\
poky/meta-poky & Poky reference distribution \\
poky/meta-yocto-bsp & Yocto BSP layer \\
meta-raspberrypi & Raspberry Pi 5 BSP \\
meta-openembedded/meta-oe & OpenEmbedded utilities \\
meta-openembedded/meta-python & Python packages \\
meta-openembedded/meta-networking & Networking packages \\
meta-openembedded/meta-multimedia & Multimedia packages \\
meta-ros/meta-ros-common & ROS common components \\
meta-ros/meta-ros2 & ROS2 core \\
meta-ros/meta-ros2-jazzy & ROS2 Jazzy distribution \\
meta-star & STAR custom recipes \\
\hline
\end{tabular}
    \caption { Yocto layers used in the STAR build}
    \end{table}

    \subsection{meta - star Layer}

        The \texttt{meta - star} layer contains STAR -
        specific recipes :

\subsubsection{recipes - core
   }

    \begin{itemize}
    \item \texttt { network - config} -- Static IP configuration for eth0 (192.168.100.2/24)
    \item \texttt{
      wifi - setup
   } -- WiFi configuration script for easy network setup
    \item \texttt{
      locale - config
   }
    --UTF - 8 locale configuration with automatic generation at boot
    \item \texttt {
      modules - autoload
   }
    --Kernel module auto - loading(i2c - dev)
    \item \texttt { ros2 - env}
    --ROS2 environment auto - sourcing
\end{itemize}

    \subsubsection{recipes - devtools}

    \begin{itemize}
    \item \texttt { temurin - jdk - bin} -- Eclipse Temurin JDK 21 for Java/Kotlin development
\end{itemize}

    \subsection{Hardware Configuration}

    The following hardware interfaces are enabled via \texttt{config.txt} :

\begin{table}[h]
\centering
\begin{tabular} {
      | l | l | l |
   }
    \hline
\textbf{Interface} & \textbf{Setting} & \textbf{Device}
    \\
\hline UART &dtoverlay = uart0 -
                          pi5 & / dev / ttyAMA0 \\
I2C &dtparam = i2c\_arm = on,
             i2c1 = on & / dev / i2c -
                    1 \\
SPI &dtparam = spi = on & / dev / spidev0 .0,
             0.1 \\
GPIO &pinctrl - rp1 &gpiochip0(54 lines) \\
Camera
                                       / Display &camera\_auto\_detect = 1,
             display\_auto\_detect =
                 1 &Auto \\
USB Power &usb\_max\_current\_enable = 1 & 5A mode \\
\hline
\end{tabular}
    \caption { Hardware interface configuration}
    \end{table}

    \subsection{Installed Packages}

    The image includes the following key packages :

\subsubsection{ROS2 Jazzy
   }
    \begin{itemize}
    \item ros - core, rclcpp, rclpy
    \item std - msgs, sensor - msgs, geometry - msgs, nav - msgs
\end{itemize}

    \subsubsection{Developer Tools}
    \begin{itemize}
    \item vim 9.1(default editor)
    \item git 2.44.4
    \item htop 3.3.0
    \item tmux 3.3a
    \item Go 1.22.12
    \item Eclipse Temurin JDK 21.0.5
\end{itemize}

    \subsubsection{Hardware Tools}
    \begin{itemize}
    \item i2c - tools(i2cdetect, i2cget, i2cset)
    \item spitools(spi - pipe, spi - config)
    \item libgpiod - tools(gpiodetect, gpioinfo, gpioget, gpioset)
\end{itemize}

    \subsubsection{Computer Vision}
    \begin{itemize}
    \item OpenCV
    \item Python3 with NumPy
    \item v4l - utils
\end{itemize}

    \subsection{Building the Image}

    \subsubsection{Prerequisites}

        Install required packages on Ubuntu :
\begin{verbatim} sudo apt install gawk wget git diffstat unzip texinfo gcc
            build -
        essential chrpath socat cpio python3 python3 - pip python3 -
        pexpect xz - utils debianutils iputils - ping python3 - git python3 -
        jinja2 python3 - subunit zstd liblz4 -
        tool file locales libacl1 bmaptool
\end{verbatim}

    \subsubsection{Setup and Build}

    \begin{verbatim} cd / path / to / star - yocto source poky / oe - init -
        build - env build -
        rpi5

#Copy configuration from star - yocto - config
                cp../
            star -
        yocto - config / conf/* conf/
ln -s ../star-yocto-config/meta-star .

bitbake core-image-minimal
\end{verbatim}

\subsubsection{Flashing}

Use bmaptool for fast, verified flashing:
\begin{verbatim}
sudo bmaptool copy\\
    tmp/deploy/images/raspberrypi5/core-image-minimal-raspberrypi5.rootfs.wic.gz\\
    /dev/sdX
\end{verbatim}

\subsection{WiFi Configuration}

The image includes a \texttt{wifi-setup} script for easy WiFi configuration:

\begin{verbatim}
# Interactive mode
wifi-setup

# Direct mode
wifi-setup <SSID> <password>

# Check status
wifi-setup --status
\end{verbatim}

\subsection{Package Management}

The image uses opkg for package management. A local package feed can be set up:

\begin{verbatim}
# On host: serve packages
cd tmp/deploy/ipk
python3 -m http.server 8080

# On Pi: configure feed
echo "src/gz all http://<host-ip>:8080/all" > /etc/opkg/feed.conf
opkg update
opkg install <package>
\end{verbatim}
